\section{Introduction}
\label{sec:intro}

Difference-in-differences (DID) is a workhorse design for evaluating climate and
environmental policy: studies of carbon pricing, emissions-trading pilots, and
emission standards compare treated and untreated units before and after a policy
to recover its effect. The credibility of any such estimate rests on a chain of
\emph{identification assumptions}---parallel trends, no anticipation, appropriate
handling of staggered adoption, no interference between units---each of which a
decade of methodological work has shown to be fragile, and which a standard
two-way fixed-effects specification can silently violate under heterogeneous,
staggered treatment timing
\cite{bertrand2004trust,goodmanbacon2021,callaway2021,sunabraham2021,dechaisemartin2020,roth2023}.

In applied work, however, these assumptions are more often \emph{asserted} than
\emph{audited}: a paper states that parallel trends ``hold'', shows an event-study
plot, and moves on. Checking whether the identification argument is actually
supported by the reported evidence is slow, expert-intensive, and inconsistent
across reviewers---one facet of a broader reliability gap in empirical research,
where flexible analytic choices and specification searching move published
conclusions \cite{silberzahn2018many,brodeur2020methods,botviniknezer2020variability}.
Large language models are an obvious candidate for scaling such checks, but they
reproduce causal \emph{language} more readily than causal \emph{inference}
\cite{kiciman2023causal,zecevic2023causalparrots,jin2024corr2cause}, so a system
that issued causal verdicts would be both unreliable and unaccountable.

\paragraph{Our approach.}
We argue that identification \emph{credibility} is auditable even when the causal
\emph{effect} is not. Because the counterfactual is never observed, we do not ask
whether an estimate is true; we ask whether each identifying assumption is
adequately supported by the evidence the paper reports, and where that support is
weakest. This turns an unanswerable causal question into a tractable auditing
problem over assumptions, their testable implications, and the in-paper evidence
that bears on them. We introduce \sys{}, a \emph{bounded-agentic} auditor: model
agency is confined to two scoped, traceable sub-tasks behind a fixed,
deterministic pipeline that decomposes a DID study into an 11-dimension rubric,
retrieves dimension-specific evidence, judges its adequacy, and localizes risk in
an evidence-anchored report for a human expert (Figure~\ref{fig:overview}). This
bounded design keeps the audit traceable and reproducible, unlike free-roaming
agents whose control flow is model-decided \cite{yao2023react}. \sys{} flags
dimensions for review; it does not adjudicate causal truth. Because the true
effect is unavailable as ground truth, we validate the auditor by \emph{flaw
injection}---perturbing a credible paper with one known flaw and checking that
\sys{} catches it, stays quiet on the clean version, and localizes the right
dimension---which needs no access to the effect. A pilot expert gold then both
evaluates and calibrates severity, since credibility is a graded, partly
design-dependent judgement rather than a binary label
\cite{jacobs2021measurement,raji2021benchmark}.

\begin{figure*}[t]
  \centering
  \includegraphics[width=\textwidth]{figures/figure1}
  \caption{Overview of \sys{}. Step~1 converts an empirical DID paper into an
  auditable structure using an 11-dimension identification rubric. Step~2
  retrieves dimension-specific evidence, assesses reported-evidence adequacy with
  a bounded LLM audit module, and localizes risk in a human-reviewable report.
  Step~3 uses expert annotation and adjudicated gold labels to evaluate and
  calibrate audit severity. \sys{} flags dimensions for expert review; it does not
  adjudicate causal truth.}
  \label{fig:overview}
\end{figure*}

\paragraph{Contributions.}
\textbf{(i)}~an \textbf{11-dimension audit framework} for DID identification
credibility whose unit of audit is an explicit
\emph{assumption--implication--evidence} chain, targeting reported-evidence
adequacy rather than causal truth (\S\ref{sec:framework});
\textbf{(ii)}~a \textbf{bounded-agentic} evidence-retrieval and assessment
pipeline, with model agency confined to two scoped, traceable sub-tasks behind a
fixed deterministic control flow;
\textbf{(iii)}~a \textbf{multi-stage evaluation} built on \emph{flaw injection as
local ground truth}, showing that adequacy reasoning catches flawed-but-present
evidence keyword matching cannot, and a real-paper deployment (27 DID papers) that
relocates the bottleneck to \emph{evidence grounding}; an architecture ablation
against a same-model single-pass baseline isolates the bounded design as the
source of the gain (\S\ref{sec:results}); and
\textbf{(iv)}~a \textbf{pilot expert-gold calibration loop} (5 papers, 55
dimension-level cells) that exposes a one-directional over-severity bias and a
lightweight, applicability-aware calibration layer that reduces it while
preserving high-risk recall.
